\section{Retain caller state and select exact memory intervals}
Preserve all sixty-four bits of \texttt{s0} through \texttt{s7}, \texttt{fp} and
\texttt{sp}, and return through the incoming \texttt{ra}. Nested calls need private
storage for their return addresses and live arguments. No routine uses console I/O.
Arguments, temporary registers and vector/overlapping floating-point storage may change.
These rules match the earlier processing kernels; a file wrapper does not relax them.

Descriptors are eight-byte aligned. PCM, coefficient, byte-file and path pointers may
be unaligned. Caller objects are disjoint. The filter reads exactly the selected PCM
and sixteen coefficient bytes and writes exactly the corresponding output PCM; the
encoder reads exactly its selected PCM and writes only its returned file bytes. The
filter and encoder read only the description prefix's first sixteen bytes. Private
stack storage and submission-owned constants do not enlarge a caller's selected slice.

For parsing and saving, reject a count above 266240 with status two and one below
forty-four with status one before pointer access. The filter validates its prefix
before PCM access or output writes. With zero frames it accesses neither PCM nor
coefficients, even though the prefix still needs validation. The bound routine also
needs a readable prefix; its interface supplies a pointer rather than packed values.
Keep these differences visible when constructing invalid-argument checks.
